\newpage
\homeworktitle{Решения домашней работы}{Версия для учителя}

\begin{teacheranswer}[Задание 1]
  \[
    \text{а) }6;
    \qquad
    \text{б) }-3;
    \qquad
    \text{в) }-2;
    \qquad
    \text{г) }\frac14.
  \]
  В пункте г): \(81^{\frac14}=3\), так как \(81=3^4\).
\end{teacheranswer}

\begin{criterionbox}
  Проверяется готовность работать не только с натуральными, но и с
  отрицательными и дробными показателями. Ошибка в пункте г) указывает не
  на непонимание логарифма, а на пробел в свойствах степеней.
\end{criterionbox}

\begin{teacheranswer}[Задание 2]
  \[
    \text{а) }\log_3 81=4;
    \qquad
    \text{б) }\log_7 1=0;
  \]
  \[
    \text{в) }\log_5\frac1{25}=-2;
    \qquad
    \text{г) }\log_{16}8=\frac34.
  \]
\end{teacheranswer}

\begin{criterionbox}
  Ключевой диагностический признак --- основание степени не меняет своей
  роли. Если ученик пишет \(\log_{81}3=4\), попросите восстановить
  проверочное равенство.
\end{criterionbox}

\begin{teacheranswer}[Задание 3]
  \[
    \text{а) }2^7=128;
    \qquad
    \text{б) }9^0=1;
  \]
  \[
    \text{в) }3^{-4}=\frac1{81};
    \qquad
    \text{г) }25^{\frac32}=125.
  \]
  Для последнего равенства:
  \[
    25^{\frac32}=\left(\sqrt{25}\right)^3=5^3=125.
  \]
\end{teacheranswer}

\section*{Решение задания 4}

\begin{teacheranswer}[Вычисления]
  \[
    \text{а) }\log_2 64=6,
    \qquad
    \text{б) }\log_7 1=0,
    \qquad
    \text{в) }\log_5\frac1{125}=-3,
  \]
  \[
    \text{г) }\log_{16}2=\frac14,
    \qquad
    \text{д) }\log_{\frac13}27=-3,
    \qquad
    \text{е) }\lg0{,}0001=-4.
  \]

  Пункт г):
  \[
    16^{?}=2,
    \qquad
    2=16^{\frac14},
    \qquad
    \log_{16}2=\frac14.
  \]

  Пункт д):
  \[
    \left(\frac13\right)^{?}=27,
    \qquad
    27=\left(\frac13\right)^{-3},
    \qquad
    \log_{\frac13}27=-3.
  \]
\end{teacheranswer}

\begin{criterionbox}
  Пункты г) и д) различают два возможных затруднения: дробный показатель и
  основание меньше единицы. Требуйте степенную проверку, а не только ответ.
\end{criterionbox}

\begin{teacheranswer}[Задание 5]
  \[
    \text{а) }2^x=32=2^5\Rightarrow \boxed{x=5};
  \]
  \[
    \text{б) }3^x=7\Rightarrow \boxed{x=\log_3 7};
  \]
  \[
    \text{в) }10^x=0{,}2\Rightarrow \boxed{x=\lg0{,}2}.
  \]
  \[
    \text{г) }4^x=\frac18
    \Longleftrightarrow
    2^{2x}=2^{-3}
    \Longleftrightarrow
    \boxed{x=-\frac32}.
  \]
\end{teacheranswer}

\begin{criterionbox}
  Проверяется понимание того, что логарифм может быть точной формой ответа.
  Если ученик записывает приближённое значение вместо \(\log_3 7\),
  проверьте, различает ли он точный и приближённый ответы.
\end{criterionbox}

\begin{teacheranswer}[Задание 6]
  \textbf{а)}
  \[
    \log_3 x=4
    \Longleftrightarrow
    x=3^4=81.
  \]

  \textbf{б)}
  \[
    \log_2(x-5)=3
    \Longleftrightarrow
    x-5=2^3=8
    \Longleftrightarrow
    \boxed{x=13}.
  \]

  \textbf{в)}
  \[
    \log_5(2x+1)=2
    \Longleftrightarrow
    2x+1=5^2=25
    \Longleftrightarrow
    \boxed{x=12}.
  \]

  \textbf{г)}
  \[
    \log_2(x-1)=-2
    \Longleftrightarrow
    x-1=2^{-2}=\frac14
    \Longleftrightarrow
    \boxed{x=\frac54}.
  \]

  Во всех пунктах аргументы логарифмов при найденных значениях положительны:
  \[
    81>0,\qquad 13-5=8>0,\qquad
    2\cdot12+1=25>0,\qquad \frac54-1=\frac14>0.
  \]
\end{teacheranswer}

\begin{criterionbox}
  Здесь проверяется перевод определения в рабочее действие. Систематическое
  изучение ОДЗ ещё не является целью комплекта, но после решения ученик
  должен проверить положительность аргумента логарифма.
\end{criterionbox}

\section*{Задания на понимание}

\begin{teacheranswer}[Задание 7]
  Имеют смысл записи
  \[
    \log_2 8\qquad\text{и}\qquad \log_{\frac12}4.
  \]
  В записи \(\log_{-2}8\) основание отрицательно, в \(\log_1 5\)
  основание равно единице, а в \(\log_3(-9)\) аргумент отрицателен.

  Далее
  \[
    \log_a16=4
    \Longleftrightarrow
    a^4=16.
  \]
  Уравнение \(a^4=16\) имеет корни \(a=2\) и \(a=-2\), но основание
  логарифма должно быть положительным и не равным единице. Поэтому
  \[
    \boxed{a=2}.
  \]
\end{teacheranswer}

\begin{criterionbox}
  Проверяется, учитывает ли ученик оба ограничения на основание
  \(a>0\), \(a\ne1\), а также положительность аргумента логарифма.
\end{criterionbox}

\begin{teacheranswer}[Задание 8]
  Логарифм отвечает не на вопрос об умножении, а на вопрос о показателе:
  «В какую степень нужно возвести 8, чтобы получить 64?»
  \[
    8^2=64,
    \qquad
    \boxed{\log_8 64=2}.
  \]
\end{teacheranswer}

\begin{criterionbox}
  Формулировка вопроса важнее исправленного числа. Она показывает, что
  ученик восстановил смысл операции, а не просто вспомнил ответ.
\end{criterionbox}

\begin{teacheranswer}[Задание 9]
  Первая запись проверяется равенством
  \[
    4^3=64.
  \]
  Вторая запись проверяется равенством
  \[
    64^{\frac13}=4.
  \]
  При смене основания изменяется вопрос: теперь нужно узнать, в какую
  степень возвести 64, чтобы получить 4. Поэтому показатель становится
  равным \(\frac13\), а не 3.
\end{teacheranswer}

\begin{criterionbox}
  Задание предупреждает механическую перестановку основания и аргумента.
  Сильный ответ содержит обе степенные проверки и словесное объяснение
  изменения вопроса.
\end{criterionbox}

\begin{teacheranswer}[Задание 10*]
  \textbf{а)} Пусть
  \[
    (\sqrt2)^x=8.
  \]
  Тогда
  \[
    \left(2^{\frac12}\right)^x=2^3
    \Longrightarrow
    \frac{x}{2}=3
    \Longrightarrow
    \boxed{x=6}.
  \]

  \textbf{б)}
  \[
    \left(\frac1{27}\right)^x=9
    \Longleftrightarrow
    (3^{-3})^x=3^2
    \Longleftrightarrow
    -3x=2
    \Longleftrightarrow
    \boxed{x=-\frac23}.
  \]

  \textbf{в)}
  \[
    (0{,}01)^x=10
    \Longleftrightarrow
    (10^{-2})^x=10^1
    \Longleftrightarrow
    -2x=1
    \Longleftrightarrow
    \boxed{x=-\frac12}.
  \]
\end{teacheranswer}

\begin{criterionbox}[Методический смысл задания со звёздочкой]
  Проверяется перенос определения в ситуацию, где ответ нельзя получить по
  непосредственному подбору. Ученик должен представить основание и аргумент
  как степени одного и того же числа и сравнить показатели.
\end{criterionbox}

\begin{criterionbox}[Общая оценка работы]
  Максимум --- 10 баллов: смысл и перевод записей (задания 1--3) ---
  3 балла; вычисление логарифмов (задание 4) --- 2 балла; уравнения и
  проверка аргументов (задания 5--6) --- 2 балла; условия существования
  логарифма (задание 7) --- 1 балл; объяснение и перенос способа
  (задания 8--9) --- 2 балла. Задание 10* можно использовать как
  дополнительное и не учитывать при снижении основной оценки.
\end{criterionbox}

\begin{resultbox}[Краткие ответы]
  \begin{enumerate}[itemsep=0.28em,topsep=0.25em]
    \item \(6;\ -3;\ -2;\ \dfrac14\).
    \item \(\log_3 81=4;\ \log_7 1=0;\
      \log_5\dfrac1{25}=-2;\ \log_{16}8=\dfrac34\).
    \item \(2^7=128;\ 9^0=1;\ 3^{-4}=\dfrac1{81};\
      25^{\frac32}=125\).
    \item \(6;\ 0;\ -3;\ \dfrac14;\ -3;\ -4\).
    \item \(5;\ \log_3 7;\ \lg0{,}2;\ -\dfrac32\).
    \item \(81;\ 13;\ 12;\ \dfrac54\).
    \item Имеют смысл \(\log_2 8\) и \(\log_{\frac12}4\); \(a=2\).
    \item \(2\).
    \item Обе записи верны.
    \item \(6;\ -\dfrac23;\ -\dfrac12\).
  \end{enumerate}
\end{resultbox}
